\section{Консольный интерфейс приложения}
Приложение решает самостоятельную пользовательскую задачу: находит самый длинный палиндром и число всех палиндромных вхождений. Для обычного расчёта пользователь не должен заранее знать ответ. Поэтому интерфейс расчёта отделён от тестового загрузчика с эталонными ожиданиями.

Основной запуск из каталога \code{benchmark_app}:
\begin{lstlisting}
npm start -- --text "abba"
npm start -- --file data/input-example.json
node src/cli.js --json '{"input":"abba"}'
node src/cli.js --text "abba" --output-json --radii
\end{lstlisting}
Последняя команда выводит машиночитаемый JSON с результатами. Имена ключей \code{longest}, \code{length}, \code{count} и \code{radii} сохраняются стабильными для автоматической обработки. Текстовые сообщения и названия тестов написаны на русском. Служебный вывод самих сторонних раннеров, например PASS и Tests, формируется их библиотеками.

Для строки \code{abba} консоль показывает:
\begin{Verbatim}[fontsize=\small,frame=single]
Самый длинный палиндром: "abba"
Длина: 4
Число палиндромных вхождений: 6
\end{Verbatim}
Одинаковые буквы на разных позициях считаются разными вхождениями. Параметр \code{--radii} добавляет массив радиусов и помогает сопоставить запуск с пошаговой таблицей алгоритма.

Если запустить \code{npm start} в интерактивном терминале без аргументов, программа запросит строку. Для пустого ввода достаточно нажать Enter. Параметр \code{--text} позволяет передать строку, начинающуюся с двух дефисов, не принимая её за параметр. Справка доступна через \code{--help}.

Пользовательский JSON допускает строку, объект с полем \code{input}, массив строк или объектов, а также объект с массивом \code{cases}. Например:
\begin{Verbatim}[fontsize=\small,frame=single]
{"cases": [
  {"name": "Чётный палиндром", "input": "abba"},
  {"name": "Русская строка", "input": "топот"},
  {"name": "Пустая строка", "input": ""}
]}
\end{Verbatim}
Поля ожидаемого ответа не требуются. При передаче тестового файла приложению они не участвуют в расчёте: проверку ожиданий выполняют команды тестов и \code{verify}. Неправильный тип \code{input}, пустой пакет, повреждённый JSON, отсутствующий файл или конфликт источников ввода дают русское сообщение в стандартном потоке ошибок и код завершения 1.

Для автоматизации предусмотрен \code{--stdin}. Совместно с \code{--json} он читает весь поток как JSON, иначе как исходную строку. Пробелы и завершающие переводы строки сохраняются. Это важно при конвейерном вводе: перевод строки, добавленный оболочкой, является частью задачи. Входной UTF-8 BOM удаляется только при разборе JSON.

Интерфейс обрабатывает Unicode по кодовым точкам без нормализации и без изменения регистра. Поэтому \code{А} и \code{а} различаются, а буква с отдельным комбинируемым знаком не равна одному предсоставленному символу. Приложение загружает JSON целиком: память зависит от суммарного объёма пакета, хотя один вызов Манакера использует линейную дополнительную память относительно длины одной строки.
